% =========================================================
% Inference Rules for Quantifiers
% =========================================================

\section{Proof Theory: Quantifier Inference Rules}

% ---------------------------------------------------------
% TOOLKIT
% ---------------------------------------------------------

% ---------------------------------------------------------
% Universal Instantiation
% ---------------------------------------------------------
\begin{definitionbox}{Definition (Universal Instantiation — UI)}
\begin{definition}[Universal Instantiation — UI]\label{def:ui}
From a universally quantified statement, infer any instance obtained by
substituting a term for the bound variable:
\[
\forall x\,\varphi \;\Rightarrow\; \varphi[t/x].
\]
\end{definition}
\end{definitionbox}

\begin{dependencies}
\begin{itemize}
  \item TODO: determine dependencies for \texttt{def:ui}.
\end{itemize}
\end{dependencies}

\begin{remark*}[Side condition]
The term $t$ must be free for $x$ in $\varphi$. If $t$ contains a variable
that would become bound after substitution, the inference is invalid.
\end{remark*}

% ---------------------------------------------------------
% Universal Generalization
% ---------------------------------------------------------
\begin{definitionbox}{Definition (Universal Generalization — UG)}
\begin{definition}[Universal Generalization — UG]\label{def:ug}
From a formula, infer a universally quantified statement:
\[
\varphi \;\Rightarrow\; \forall x\,\varphi.
\]
\end{definition}
\end{definitionbox}

\begin{dependencies}
\begin{itemize}
  \item TODO: determine dependencies for \texttt{def:ug}.
\end{itemize}
\end{dependencies}

\begin{remark*}[Restriction on UG]
The variable $x$ must not occur free in any undischarged assumption on which
$\varphi$ depends. This restriction ensures that $x$ is truly arbitrary —
not tied to a specific hypothesis about $x$.
\end{remark*}

\begin{remark*}[Arbitrary element argument]
To prove $\forall x\,\varphi(x)$, fix an arbitrary element $x$ (introducing no
special assumptions about $x$), prove $\varphi(x)$, and then apply UG. The
variable $x$ must remain unconstrained throughout.
\end{remark*}

% ---------------------------------------------------------
% Existential Instantiation
% ---------------------------------------------------------
\begin{definitionbox}{Definition (Existential Instantiation — EI)}
\begin{definition}[Existential Instantiation — EI]\label{def:ei}
From an existentially quantified statement, infer an instance with a fresh
constant (witness):
\[
\exists x\,\varphi \;\Rightarrow\; \varphi[c/x].
\]
\end{definition}
\end{definitionbox}

\begin{dependencies}
\begin{itemize}
  \item TODO: determine dependencies for \texttt{def:ei}.
\end{itemize}
\end{dependencies}

\begin{remark*}[Witness discipline for EI]
The constant $c$ must be fresh: it must not appear in $\varphi$, in any
undischarged assumption, or in the final conclusion of the proof. The constant
represents an arbitrary but fixed witness to the existential claim, not a
specific known object.
\end{remark*}

\begin{remark*}[Common error]
Using a witness constant introduced by EI outside its allowed scope, or
choosing $c$ before establishing the existential premise, invalidates the proof.
\end{remark*}

% ---------------------------------------------------------
% Existential Generalization
% ---------------------------------------------------------
\begin{definitionbox}{Definition (Existential Generalization — EG)}
\begin{definition}[Existential Generalization — EG]\label{def:eg}
From a formula containing a term, infer an existential statement:
\[
\varphi[t/x] \;\Rightarrow\; \exists x\,\varphi.
\]
\end{definition}
\end{definitionbox}

\begin{dependencies}
\begin{itemize}
  \item TODO: determine dependencies for \texttt{def:eg}.
\end{itemize}
\end{dependencies}

\begin{remark*}[No side conditions]
EG has no side conditions: any term $t$ may be replaced by an existentially
quantified variable. To prove $\exists x\,\varphi(x)$, explicitly exhibit a
term $t$ such that $\varphi(t)$ holds, then apply EG.
\end{remark*}

\begin{remark*}[Counterexample argument]
To refute $\forall x\,\varphi(x)$: produce a single term $t$ such that
$\neg\varphi(t)$ holds. To refute $\exists x\,\varphi(x)$: show that $\varphi(t)$
fails for every term $t$ in the domain.
\end{remark*}
